% intro：生成于阅读卡+证据；引用检查通过
\section{引言}

大语言模型（Large Language Model, LLM）驱动的智能体在长期交互任务中面临一个核心约束：模型的原生上下文窗口是有限且易失的。现有研究表明，即便上下文窗口被扩展，它仅提供会话内的召回能力，并不能自动实现跨会话的持久保留、共享与检索，因此不构成真正意义上的长期记忆（long-term memory, LTM）\cite{W0039}。基于提示的 LLM 智能体受上下文窗口大小限制，无法回忆已见内容，也难以在少量演示之外持续学习\cite{W0618}。在静态检索增强生成（Retrieval-Augmented Generation, RAG）范式下，检索通常针对单一查询执行，缺乏面向跨会话可写状态的读写接口；而参数化知识既难以在推理时被精确更新，也难以携带个体经验与交互轨迹。作者在 MemGPT 中进一步指出，LLM 需要借助外部存储与类似操作系统的自主记忆管理，才能突破固定上下文窗口的约束\cite{W0002}。这些限制共同构成了长期任务引入外部可写记忆的基本动因。

为此，研究者已提出多种长期记忆实现方案。生成式智能体（Generative Agents）以自然语言记忆流记录智能体全部经验，通过相关性、新近性与重要性综合筛选，并引入反思机制将记忆综合为高层推断\cite{W0001}。MemGPT 将主上下文划分为只读系统指令、可写工作上下文与消息队列，外部上下文类比磁盘存储，记忆编辑与检索完全由 LLM 自主完成\cite{W0002}。RET-LLM 旨在提供通用读写记忆，支持跨任务知识存储与检索，且不修改 LLM 架构\cite{W0680}。ExpeL 通过自主收集经验并提取跨任务洞察实现学习与回忆\cite{W0618}。Mem0 与 A-Mem 则分别从生产级记忆操作与动态组织机制角度，尝试构建可扩展、可演化的记忆系统\cite{W0682,W0681}。然而，作者指出当前系统多停留在基础存储与检索层面，即使引入图数据库，也可能因固定操作与预定 schema 而受限\cite{W0681}。同时，术语与边界的不统一增加了横向比较的困难；CoALA 框架从认知架构角度区分了程序性、语义与情景三类记忆，为理解智能体内部状态提供了概念参照\cite{W0140}。安全综述则将长期记忆界定为可写、跨会话的持久记忆，并提出涵盖写入、存储、检索、执行、共享与传播、遗忘与回滚的生命周期框架\cite{W0039}。

在上述背景下，本文关注长期任务中外部记忆机制的读写闭环及其对智能体行为的影响，而非参数化记忆更新，也不研究单查询检索式 RAG 的检索质量本身。围绕跨会话可写外部记忆，本文提出以下研究问题：

\textbf{RQ1（定义与边界）：}长期记忆包含哪些信息和机制？它与上下文历史、静态知识检索、参数知识有什么关系？

\textbf{RQ2（组织与使用）：}写入、组织和检索方法如何影响信息可用性、任务表现和成本？

\textbf{RQ3（更新与遗忘）：}如何处理信息变化、时间条件、冲突、过时记忆和必要的删除？

\textbf{RQ4（可靠性）：}错误摘要、错误反思或来源混淆如何进入记忆并影响后续决策？现有方法怎样发现和控制这些问题？

\textbf{RQ5（评测）：}基准和实验是否真正考察跨会话记忆、更新能力和长期任务表现？不同结果能否公平比较？

上述研究问题限定在跨会话可写的外部记忆范围之内。本文综合已有工作，但不将参数化记忆纳入讨论，也不对单查询检索式 RAG 的索引与排序算法作系统比较。对于跨系统、跨任务的可比较度量与结论，若素材未直接给出，本文仅作定性综合，并明确标注为【待核验】。
